% Elaborado por Fabian Gomez
% Version 1.0

% =========================================================
% ICD LLC — Interfaz Arduino Mega 2560 ↔ Raspberry Pi 5
% =========================================================
\subsection{ICD LLC (Interfaz Arduino Mega\,2560\,R3 -- Raspberry Pi\,5)}
\label{sec:icd_llc}

\subsubsection{Alcance}
Este ICD define la interfaz \textbf{interna} serie entre el Nodo~B
(MCU -- Arduino Mega\,2560\,R3, bajo nivel) y el Nodo~A
(OBC -- Raspberry Pi\,5, alto nivel) del Rover Olympus.
Especifica el medio físico, el protocolo ASCII delimitado por
\verb|\n| y el diccionario completo de mensajes
(comandos CMD y respuestas RSP) tal como están implementados
en el firmware LLC (Rust, v2.8).

\textbf{Mapeo con subsistemas SyRS:}
\begin{itemize}
  \item \textbf{Nodo A (RPi\,5)} = C\&DH/OBC; aloja además la lógica GNC
        (ELANav: percepción, planificación, navegación) y la capa COMMS (CSP/HLC).
  \item \textbf{Nodo B (Arduino Mega)} = control de bajo nivel para PROP
        (motores PWM), TRAC (encoders Hall) y adquisición de sensores
        (HC-SR04, VL53L0X, INA3221, NTC).
  \item \textbf{Fuera de scope de este ICD:} interfaz CSI-2 de la cámara
        OV5647 (conectada directamente a RPi\,5); lógica de percepción/navegación
        (GNC); stack COMMS (CSP/UHF/Wi-Fi). Ver ICD-CAM-001 e ICD-COMMS-001.
\end{itemize}

Las responsabilidades y requisitos de sistema a los que este ICD da
cumplimiento son: SyRS~§7.4 (UART), §8 (arquitectura Nodo\,A/B),
PROP\-REQ\-002 (parada de motores $\leq 500$\,ms),
RF‑005 (Fault Stop) y RNF‑001 ($\leq 2$\,s por ciclo de agente).

% ---------------------------------------------------------
\subsubsection{Entidades y roles}
\begin{itemize}
  \item \textbf{RPi\,5 (Maestro / Iniciador):}
        genera todos los comandos CMD; lee y procesa las respuestas RSP.
  \item \textbf{Arduino Mega (Esclavo / Ejecutor):}
        recibe CMD, ejecuta la transición de estado y responde
        con RSP o TLM.
\end{itemize}

% ---------------------------------------------------------
\subsubsection{Interfaz física}

\paragraph{IF-LLC-UART (Arduino Mega $\leftrightarrow$ RPi\,5).}
\begin{itemize}
  \item \textbf{Puerto Arduino:} USART0 → pines D0 (RX) / D1 (TX) del Mega\,2560.
  \item \textbf{Puerto RPi\,5:} \texttt{/dev/ttyUSB0} o \texttt{/dev/ttyACM0}
        según cable USB–Serie empleado.
  \item \textbf{Velocidad:} 115\,200 baud.
  \item \textbf{Formato:} 8N1 (8 bits de datos, sin paridad, 1 bit de parada).
  \item \textbf{Control de flujo:} ninguno (sin RTS/CTS).
  \item \textbf{Nivel lógico:} 5\,V TTL en los pines del Mega;
        el adaptador USB–Serie realiza la conversión a los niveles del host.
\end{itemize}
(\emph{Ref. SyRS §7.3 hardware interfaces, ASM-042})

% ---------------------------------------------------------
\subsubsection{Protocolo de comunicación}

Todos los mensajes son cadenas ASCII terminadas en \verb|\n| (0x0A).
Se tolera \verb|\r\n| en la recepción; el emisor \textbf{siempre} termina
con \verb|\n| solamente.

\paragraph{Modelo transaccional.}
La RPi\,5 envía un CMD y el Arduino responde con exactamente una RSP
en el mismo ciclo de loop ($\leq 20$\,ms).
La telemetría extendida (TLM) se emite \textbf{adicionalmente} de forma
periódica sin esperar un CMD específico.

\paragraph{Límites de tamaño.}
\begin{center}
\small
\begin{tabular}{lll}
\toprule
\textbf{Buffer} & \textbf{Tamaño} & \textbf{Propósito} \\
\midrule
CMD (Arduino RX) & 32 bytes & Frame de comando desde RPi\,5 \\
RSP (Arduino TX) & 160 bytes & Frame de respuesta / TLM extendido \\
Ring RX (ISR)    & 64 bytes & Buffer circular llenado por ISR USART\_RX \\
\bottomrule
\end{tabular}
\end{center}

% ---------------------------------------------------------
\subsubsection{Diccionario de comandos (RPi\,5 $\rightarrow$ Arduino)}

\begin{center}
\small
\renewcommand{\arraystretch}{1.3}
\begin{tabular}{p{3.4cm}p{4.5cm}p{5.5cm}}
\toprule
\textbf{Comando} & \textbf{Descripción} & \textbf{Notas} \\
\midrule
\verb|PING\n| &
  Keepalive / watchdog reset &
  Respuesta: \verb|PONG|. Debe enviarse al menos cada $\leq 2$\,s
  en estados activos o el Arduino entra en FAULT. \\

\verb|STB\n| &
  Transición a Standby &
  Detiene todos los motores. Válido en cualquier estado salvo FAULT. \\

\verb|EXP:<L>:<R>\n| &
  Comando de exploración &
  \verb|<L>| y \verb|<R>| son enteros con signo $[-99, +99]$
  (porcentaje de potencia). Ejemplo: \verb|EXP:80:80\n|.\\

\verb|AVD:L\n| / \verb|AVD:R\n| &
  Maniobra de evasión (izq. / der.) &
  Aplica $\pm 60$\,\% de potencia en skid-steer. \\

\verb|RET\n| &
  Retroceso &
  Aplica $-50$\,\% de potencia en ambos lados. \\

\verb|FLT\n| &
  Forzar estado FAULT &
  Para motores de inmediato. Solo \verb|RST| sale de FAULT. \\

\verb|RST\n| &
  Reset del sistema &
  Vuelve a Standby desde cualquier estado, incluido FAULT.
  Borra safety level y watchdog. \\
\bottomrule
\end{tabular}
\end{center}
(\emph{Ref. firmware LLC \texttt{state\_machine/mod.rs}: \texttt{parse\_command()}, \texttt{Command} enum})

% ---------------------------------------------------------
\subsubsection{Diccionario de respuestas (Arduino $\rightarrow$ RPi\,5)}

\begin{center}
\small
\renewcommand{\arraystretch}{1.3}
\begin{tabular}{p{4.5cm}p{3.0cm}p{5.4cm}}
\toprule
\textbf{Respuesta} & \textbf{Disparada por} & \textbf{Descripción} \\
\midrule
\verb|PONG\n| &
  CMD \verb|PING| &
  Confirma disponibilidad y resetea el watchdog. \\

\verb|ACK:<STATE>\n| &
  CMD exitoso &
  Confirma la transición. \verb|<STATE>| $\in$ \{STB, EXP, AVD, RET, FLT\}.
  Ejemplo: \verb|ACK:EXP\n|. \\

\verb|TLM:<SAF>:<STALL>:| \newline
\verb|<TS>ms:<MV>mV:<MA>mA:| \newline
\verb|<I0>:<I1>:<I2>:<I3>:| \newline
\verb|<I4>:<I5>:<T>C:| \newline
\verb|<B0>:<B1>:<B2>:| \newline
\verb|<B3>:<B4>:<B5>C:| \newline
\verb|<DIST>mm\n| &
  Periódico ($\approx 1$\,s) &
  Frame de telemetría extendida. Ver §\ref{ssec:icd_llc_tlm}. \\

\verb|ERR:ESTOP\n| &
  CMD en estado FAULT &
  El Arduino está en FAULT; solo acepta \verb|RST|. \\

\verb|ERR:UNKNOWN\n| &
  CMD no reconocido &
  El frame recibido no coincide con ningún comando válido. \\

\verb|ERR:WDOG\n| &
  Expiración de watchdog &
  El Arduino no recibió ningún CMD en $\geq 100$ ciclos
  ($\approx 2$\,s) estando en estado activo; transiciona a FAULT. \\
\bottomrule
\end{tabular}
\end{center}
(\emph{Ref. firmware LLC \texttt{state\_machine/mod.rs}: \texttt{format\_response()}, \texttt{Response} enum})

% ---------------------------------------------------------
\subsubsection{Frame de telemetría extendida (TLM)}
\label{ssec:icd_llc_tlm}

Formato completo (máx. $\approx 140$ bytes + \verb|\n|):

\begin{center}
\small\ttfamily
TLM:<SAF>:<STALL>:<TS>ms:<MV>mV:<MA>mA:<I0>:<I1>:<I2>:<I3>:\\
\phantom{TLM:}<I4>:<I5>:<T>C:<B0>:<B1>:<B2>:<B3>:<B4>:<B5>C:<DIST>mm\textbackslash n
\end{center}

\begin{center}
\small
\renewcommand{\arraystretch}{1.3}
\begin{tabular}{p{1.8cm}p{3.0cm}p{7.5cm}}
\toprule
\textbf{Campo} & \textbf{Tipo / Rango} & \textbf{Descripción} \\
\midrule
\verb|<SAF>| &
  \{NORMAL, WARN, LIMIT, FAULT\} &
  Nivel de seguridad actual de la MSM Nodo\,B. \\

\verb|<STALL>| &
  6 bits ASCII '0'/'1' &
  Máscara de stall por motor: bit5$\ldots$bit0 $\equiv$ motor5$\ldots$motor0
  (FR, FL, CR, CL, RR, RL). Ejemplo: \verb|000001| = stall en motor RL. \\

\verb|<TS>| &
  u32, [ms] &
  Tiempo relativo desde arranque del firmware (contador monotónico, overflow a $\approx 49$ días). \\

\verb|<MV>| &
  u16, [mV] &
  Tensión del bus de batería medida por INA3221 (pines D42/D43 soft-I2C).
  \textbf{0} = sensor no disponible. \\

\verb|<MA>| &
  i32, [mA] &
  Corriente total de batería con signo medida por INA3221.
  \textbf{0} = sensor no disponible. \\

\verb|<I0>–<I5>| &
  i32 $\times$ 6, [mA] &
  Corriente por motor [FR, FL, CR, CL, RR, RL] medida por ACS712
  (pines A0–A5). Puede ser negativa. \\

\verb|<T>| &
  i32, [°C] &
  Temperatura ambiente medida por LM335 (pin A6). \\

\verb|<B0>–<B5>| &
  i32 $\times$ 6, [°C] &
  Temperatura por sensor NTC de batería [B1a, B1b, B2a, B2b, B3a, B3b]
  (pines A7–A12). \\

\verb|<DIST>| &
  u16, [mm] &
  Distancia medida por VL53L0X ToF (soft-I2C D42/D43).
  \textbf{0} = sin lectura disponible. \\
\bottomrule
\end{tabular}
\end{center}

Ejemplo de frame TLM real:
\begin{verbatim}
TLM:NORMAL:000000:12340ms:11800mV:2350mA:
    200:210:195:205:180:190:24C:
    25:25:26:26:25:25C:450mm
\end{verbatim}
(\emph{Ref. firmware LLC \texttt{state\_machine/mod.rs}: \texttt{format\_tlm()}, \texttt{SensorFrame} struct})

% ---------------------------------------------------------
\subsubsection{Máquina de estados y comandos permitidos}

La MSM del Nodo\,B define los estados operativos y las transiciones
permitidas. Los comandos en \textbf{negrita} son los que causan
transición de estado.

\begin{center}
\small
\renewcommand{\arraystretch}{1.3}
\begin{tabular}{lp{4.5cm}p{5.0cm}}
\toprule
\textbf{Estado} & \textbf{Comandos aceptados} & \textbf{Respuesta a otros CMD} \\
\midrule
STANDBY &
  PING, \textbf{EXP}, \textbf{AVD}, \textbf{RET}, \textbf{FLT}, \textbf{RST}, STB &
  \verb|ERR:UNKNOWN| si no reconocido \\
EXPLORE &
  PING, \textbf{STB}, \textbf{AVD}, \textbf{RET}, \textbf{FLT}, \textbf{RST} &
  Watchdog activo; sin PING $\geq 2$\,s → FAULT \\
AVOID   &
  PING, \textbf{STB}, \textbf{EXP}, \textbf{RET}, \textbf{FLT}, \textbf{RST} &
  Watchdog activo \\
RETREAT &
  PING, \textbf{STB}, \textbf{EXP}, \textbf{AVD}, \textbf{FLT}, \textbf{RST} &
  Watchdog activo \\
FAULT   &
  \textbf{RST} únicamente &
  \verb|ERR:ESTOP| para cualquier otro CMD \\
\bottomrule
\end{tabular}
\end{center}

Transiciones autónomas (sin CMD externo):
\begin{itemize}
  \item \textbf{Watchdog expira} (estados activos): cualquier estado activo $\rightarrow$ FAULT,
        emite \verb|ERR:WDOG|.
  \item \textbf{Obstáculo HC-SR04} ($< 200$\,mm): cualquier estado $\rightarrow$ FAULT.
  \item \textbf{Obstáculo VL53L0X} ($< 150$\,mm): cualquier estado $\rightarrow$ FAULT.
  \item \textbf{Stall de encoder}: motor$_N$ sin pulsos durante $\geq 50$ ciclos a
        velocidad $> 20$\,\% → FAULT, stall\_mask bit$_N = 1$.
  \item \textbf{Sobrecorriente ACS712}: umbral WARN (1\,200\,mA) o
        LIMIT (1\,600\,mA) actualiza \verb|SafetyState|; umbral FAULT → FAULT.
\end{itemize}
(\emph{Ref. firmware LLC \texttt{main.rs} v2.8, \texttt{state\_machine/mod.rs}})

% ---------------------------------------------------------
\subsubsection{Temporización}

\begin{center}
\small
\renewcommand{\arraystretch}{1.3}
\begin{tabular}{lll}
\toprule
\textbf{Parámetro} & \textbf{Valor} & \textbf{Notas} \\
\midrule
Período de loop              & 20\,ms     & Constante \texttt{LOOP\_MS} \\
Watchdog timeout             & $\approx 2$\,s (100 ciclos) & Constante \texttt{WATCHDOG\_MAX} \\
Período telemetría TLM       & $\approx 1$\,s (50 ciclos)  & Constante \texttt{TLM\_PERIOD} \\
Período lectura HC-SR04      & $\approx 100$\,ms (5 ciclos) & Constante \texttt{HC\_READ\_PERIOD} \\
Período lectura sensores ADC & $\approx 500$\,ms (25 ciclos) & Constante \texttt{SEN\_READ\_PERIOD} \\
Latencia CMD $\rightarrow$ RSP & $\leq 20$\,ms & Una respuesta por ciclo \\
Motor stop tras FAULT        & $\leq 20$\,ms & Dentro del mismo ciclo \\
\bottomrule
\end{tabular}
\end{center}

\noindent\textbf{Nota sobre PROP-REQ-002:} el requisito impone
parada de motores $\leq 500$\,ms desde recepción de Fault Stop.
El Arduino cumple en $\leq 20$\,ms (un ciclo);
el tiempo restante corresponde al presupuesto de detección
y comando en el Nodo\,A (RPi\,5).

% ---------------------------------------------------------
\subsubsection{Asignación de pines (resumen)}

\begin{center}
\small
\renewcommand{\arraystretch}{1.3}
\begin{tabular}{lll}
\toprule
\textbf{Pin / Bus} & \textbf{Función} & \textbf{Interfaz} \\
\midrule
D0 (RX) / D1 (TX)   & USART0 → RPi\,5 / USB    & IF-LLC-UART \\
D5, D6, D7, D8, D9, D10 & PWM motores (Timers 2/3/4) & PROP \\
D22–D37              & GPIO DIR motores           & PROP \\
D38 (Trigger), D39 (Echo) & HC-SR04 ultrasónico  & GNC \\
D42 (SDA), D43 (SCL) & Soft I2C → VL53L0X + INA3221 & GNC / EPS \\
D18–D21, D2, D3      & INT encoder Hall (FR..RL)  & TRAC \\
A0–A5                & ACS712 corriente motores   & EPS \\
A6                   & LM335 temperatura ambiente & GNC \\
A7–A12               & NTC temperatura baterías   & EPS \\
\bottomrule
\end{tabular}
\end{center}
(\emph{Ref. firmware LLC \texttt{main.rs} v2.8 §Asignación de pines})

% ---------------------------------------------------------
\subsubsection{Mapeo PBS/ASM}

\begin{itemize}
  \item \textbf{ASM-042} (MCU Arduino Mega 2560 R3):
        [DATA] USART0 $\leftrightarrow$ RPi\,5 (IF-LLC-UART);
        [CTRL] PWM+DIR $\rightarrow$ BTS7960 (PROP);
        [DATA] soft-I2C $\rightarrow$ VL53L0X / INA3221;
        [DATA] GPIO $\rightarrow$ HC-SR04;
        [DATA] GPIO ISR $\leftarrow$ encoders Hall (TRAC);
        [DATA] ADC $\leftarrow$ ACS712 / LM335 / NTC.
\end{itemize}

% ---------------------------------------------------------
\subsubsection{Verificación de la interfaz (ICD checks)}

\begin{enumerate}
  \item \textbf{Enlace físico}: conectar RPi\,5 al Arduino Mega por USB;
        ejecutar \texttt{read\_serial.py} (incluido en el repo LLC);
        enviar \verb|PING\n| y verificar respuesta \verb|PONG\n| en $\leq 20$\,ms.
  \item \textbf{Comandos básicos}: ciclar STB, EXP, AVD, RET y verificar
        ACK correcto en cada caso.
  \item \textbf{Fault Stop (PROP-REQ-002)}: enviar \verb|EXP:80:80\n|;
        enviar \verb|FLT\n|; medir tiempo entre CMD y parada eléctrica
        de motores; criterio: $\leq 500$\,ms desde emisión del CMD.
  \item \textbf{Watchdog}: enviar \verb|EXP:50:50\n|; detener envío de PING
        $\geq 2$\,s; verificar recepción de \verb|ERR:WDOG\n| y parada de motores.
  \item \textbf{Telemetría TLM}: verificar recepción de frame TLM cada $\approx 1$\,s;
        parsear campos SAF, STALL, TS, MV, MA, DIST y validar rango.
  \item \textbf{ERR:ESTOP}: poner Arduino en FAULT mediante \verb|FLT\n|;
        enviar cualquier CMD excepto \verb|RST\n|; verificar \verb|ERR:ESTOP\n|.
\end{enumerate}
(\emph{Ref. SyRS V\&V-PROP-002; firmware LLC \texttt{read\_serial.py}})

% ---------------------------------------------------------
\subsubsection{Constantes de calibración y odometría}
\label{sssec:icd_llc_calib}

El Nodo~B acumula pulsos de encoder y reporta los contadores brutos
(\texttt{I0}–\texttt{I5}) en cada trama TLM.
La conversión a distancia lineal y orientación la realiza el Nodo~A (HLC)
usando las constantes físicas de la Tabla~\ref{tab:icd_llc_odo_const}.

\begin{table}[H]
\centering
\small
\renewcommand{\arraystretch}{1.3}
\begin{tabular}{llp{5.5cm}}
\toprule
\textbf{Constante} & \textbf{Valor por defecto} & \textbf{Procedimiento de calibración} \\
\midrule
\texttt{TICKS\_PER\_REV}  & 20 (\textbf{TBD}) &
  Elevar rover, girar una rueda una vuelta completa y contar pulsos con
  \texttt{test\_encoders}. Repetir 5 veces y promediar. \\
\texttt{WHEEL\_RADIUS\_MM} & 50\,mm (\textbf{TBD}) &
  Medir diámetro de la rueda con calibrador y dividir entre dos. \\
\texttt{WHEEL\_BASE\_MM}  & 280\,mm (\textbf{TBD}) &
  Medir separación entre centros de contacto de ruedas izquierda y
  derecha sobre terreno llano. \\
\bottomrule
\end{tabular}
\caption{Constantes de odometría pendientes de calibración en campo.}
\label{tab:icd_llc_odo_const}
\end{table}

\noindent\textbf{Nota:} los valores por defecto permiten verificar el flujo
completo LLC\,$\rightarrow$\,TLM\,$\rightarrow$\,HLC sin hardware real.
La precisión requerida por RNF-003 ($<5\%$ de error) no puede validarse
hasta completar la calibración con hardware en campo.
(\emph{Ref. RNF-003; tesis §Constantes de calibración})

% ---------------------------------------------------------
\subsubsection{Monitores de seguridad HLC}
\label{sssec:icd_llc_monitors}

El Nodo~A (HLC) implementa seis monitores que operan sobre cada trama TLM
y aplican la política de seguridad multicapa (RF-005, SyRS-080).
Cada monitor expone \texttt{update(tlm)} y retorna un nivel de severidad;
el \texttt{HlcEngine} los consulta en orden de prioridad decreciente
antes de despachar el comando de la fuente activa.

\begin{table}[H]
\centering
\small
\renewcommand{\arraystretch}{1.3}
\begin{tabular}{p{3.8cm} l p{6.5cm}}
\toprule
\textbf{Monitor} & \textbf{Requisito} & \textbf{Umbrales / lógica} \\
\midrule
\texttt{EnergyMonitor} & EPS-REQ-001 &
  \texttt{BATT\_WARN\_MV}\,=\,14\,000\,mV (3.5\,V/celda 4S);\newline
  \texttt{BATT\_CRITICAL\_MV}\,=\,12\,800\,mV (3.2\,V/celda).\newline
  \texttt{CRITICAL} activa \texttt{SafeMode}; solo \texttt{RST}
  con condición resuelta lo desactiva. Lecturas 0\,mV descartadas. \\
\texttt{ThermalMonitor} & RNF-004 &
  \texttt{TEMP\_WARN\_C}\,=\,45\,°C;\newline
  \texttt{TEMP\_CRIT\_C}\,=\,60\,°C.\newline
  Umbral crítico activa \texttt{SafeMode}. Lecturas 0\,°C descartadas. \\
\texttt{SlipMonitor} & RF-004 &
  \texttt{SLIP\_STALL\_FRAMES}\,=\,2 ciclos EXP consecutivos con
  \texttt{stall\_mask}\,$\neq$\,\texttt{000000}. Al superarlo, emite
  \texttt{RET}. Contador se reinicia al salir de EXP o con stall\_mask=0. \\
\texttt{WaypointTracker} & SYS-FUN-020 &
  Cola FIFO de hasta \texttt{MAX\_WAYPOINTS}\,=\,5 posiciones seguras
  (solo registra en EXP con \texttt{safety=NORMAL}).\newline
  \texttt{RETREAT\_DIST\_MM}\,=\,300\,mm: emite \texttt{RET} proactivo
  antes del umbral hardware LLC (200\,mm). \\
\texttt{SafeMode} & RF-005, SYS-FUN-040 &
  Se activa por: batería crítica \textbf{o} \texttt{safety=FAULT} en TLM
  \textbf{o} temperatura crítica.\newline
  Bloquea todo comando excepto \texttt{STB}, \texttt{PING}, \texttt{RST}.\newline
  Exclusivo del HLC (no confundir con watchdog hardware del LLC). \\
\texttt{CommLinkMonitor} & SyRS-051, SRS-013 &
  Estados: \textit{COMUNICAR}\,/\,\textit{GESTION\_ENLACE}.\newline
  \texttt{GCS\_LINK\_LOST\_S}\,=\,10\,s sin trama GCS válida → transición.\newline
  \texttt{GCS\_MAX\_RETRIES}\,=\,3 intentos con
  \texttt{GCS\_RETRY\_INTERVAL\_S}\,=\,5\,s.\newline
  Sin respuesta → activa \texttt{SafeMode}. \\
\bottomrule
\end{tabular}
\caption{Monitores de seguridad HLC y sus umbrales operacionales.}
\label{tab:icd_llc_monitors}
\end{table}

(\emph{Ref. RF-005, SyRS-080, EPS-REQ-001, RNF-004, RF-004; tesis §Monitores de seguridad})

% ---------------------------------------------------------
\subsubsection{Pipeline de percepción reactiva (VisionSource)}
\label{sssec:icd_llc_vision}

Durante el estado \texttt{EXP}, el Nodo~A puede operar en modo autónomo
guiado por \texttt{VisionSource}, que adquiere un fotograma CSI,
ejecuta inferencia con YOLOv8n(-seg) y mapea el resultado a un
comando LLC (\texttt{EXP}, \texttt{AVD:L}, \texttt{AVD:R} o \texttt{RET}).

\paragraph{Parámetros configurables (\texttt{olympus\_controller.yaml}).}
\begin{table}[H]
\centering
\small
\renewcommand{\arraystretch}{1.3}
\begin{tabular}{lll}
\toprule
\textbf{Parámetro} & \textbf{Valor por defecto} & \textbf{Descripción} \\
\midrule
\texttt{VISION\_MODE}      & \texttt{bbox}  & \texttt{bbox} (YOLOv8n) o \texttt{segmentation} (YOLOv8n-seg) \\
\texttt{VISION\_CONF\_MIN} & 0.50           & Confianza mínima de detección (bbox) \\
\texttt{VISION\_AREA\_MIN} & 0.05           & Fracción mínima de área de frame (bbox) \\
\texttt{SEG\_CONF\_MIN}    & 0.50           & Confianza mínima (seg) \\
\texttt{SEG\_AREA\_MIN}    & 0.03           & Fracción mínima de área (seg) \\
\texttt{SEG\_ZONE\_MIN}    & 0.05           & Cobertura mínima para emitir avoidance \\
\texttt{SEG\_ROI\_TOP}     & 0.50           & Fracción superior del frame excluida del ROI \\
\texttt{ZONE\_LEFT\_END}   & 0.33           & Límite derecho de zona izquierda (normalizado) \\
\texttt{ZONE\_RIGHT\_START} & 0.67          & Límite izquierdo de zona derecha (normalizado) \\
\texttt{EXP\_SPEED\_L/R}   & 40             & Velocidad de exploración (\%) \\
\bottomrule
\end{tabular}
\caption{Parámetros configurables de \texttt{VisionSource}.}
\label{tab:icd_llc_vision_params}
\end{table}

\paragraph{Mapeo zona\,$\rightarrow$\,comando LLC.}
\begin{table}[H]
\centering
\small
\renewcommand{\arraystretch}{1.3}
\begin{tabular}{p{4.2cm} p{2.8cm} p{2.8cm} l}
\toprule
\textbf{Condición} & \textbf{Modo bbox} & \textbf{Modo seg} & \textbf{Comando} \\
\midrule
Sin detección / cov.\,$<$\,umbral & $c_x =$ ninguno & $\text{cov}_\text{max} <$ \texttt{SEG\_ZONE\_MIN} & \texttt{EXP:40:40} \\
Detección zona izquierda          & $c_x < 0.33$    & $\text{cov}_\text{izq}$ máxima              & \texttt{AVD:R} \\
Detección zona central            & $0.33 \le c_x \le 0.67$ & $\text{cov}_\text{centro}$ máxima  & \texttt{RET} \\
Detección zona derecha            & $c_x > 0.67$    & $\text{cov}_\text{der}$ máxima              & \texttt{AVD:L} \\
\bottomrule
\end{tabular}
\caption{Mapeo zona\,$\rightarrow$\,comando MSM en \texttt{VisionSource}.}
\label{tab:icd_llc_vision_cmd}
\end{table}

\noindent\textbf{Nota:} en modo \texttt{segmentation}, si el archivo
\texttt{yolov8n-seg.onnx} no existe en la ruta configurada, el sistema
cae automáticamente a modo \texttt{bbox} con aviso en log.
(\emph{Ref. GNC-REQ-001, GNC-REQ-002, RF-001; tesis §Pipeline de visión})

% ---------------------------------------------------------
\subsubsection{Cumplimiento con SyRS}

Este ICD satisface:
SyRS~§7.3 (UART, GPIO, PWM), §7.4 (interfaces internas),
§8 (arquitectura Nodo\,A/B), ASM-042,
PROP\-REQ\-002 (motor stop $\leq 500$\,ms),
RF‑005 (Fault Stop), RNF‑001 (latencia $\leq 2$\,s por ciclo).%
